% Details moved from the main text (methods, protocol and secondary results); every number is a \val macro.
\subsection{Error structure}
The comparison of spread models in Section~\ref{sec:method} of the main paper uses \val{ENPairsMid} and
\val{ENPairsHigh} pairs of lead time and error from \val{ENMissions} development missions per disturbance level. A
hybrid model $s(h)=\sqrt{q_0+q_1h+q_2h^2}$ also fitted worse than the linear spread, by \val{EDAicHybMid} and
\val{EDAicHybHigh} points of the Akaike information criterion. The flight-time error of a single leg had a coefficient of
variation of \val{EMpCvMid} and \val{EMpCvHigh} under mid and high disturbances, and missions explained only
\val{EShareMissionMid}\% and \val{EShareMissionHigh}\% of the residual variance.

\subsection{Parameters of the compared methods}\label{sec:Sparams}
The performance profile rule predicts the gain of further search from a static profile over the elapsed time and the
number of open tasks. The online prediction rule follows Svegliato et al.~\cite{svegliato2018meta}, but fits the
saturating curve $q(t)=a-b\,e^{-ct}$ to the best objective found so far at every check, in place of the arctangent or
logarithmic models of the original. Both use the delay cost and the price of computation of AnytimeSwarm-R. The size-scaled budget allows 1\,ms
per open task, and the stagnation rule stops after 20\% of the budget without improvement. The proportional buffer
accepts a planned finish $f_j$ if $f_j+c\,h_j\le d_j$. Moment propagation models the flight time of a leg of nominal
duration $t$ as normal with mean $(1+a)t+c_0$ and standard deviation $c_vt$, and the service time with mean $(1+a_s)s$
and standard deviation $c_ss$. Legs are independent, the start of service $\max(A_j,r_j)$ is approximated by Clark's
moments~\cite{clark1961greatest}, and a finish is accepted if $\mathrm{E}[F_j]+z\,\mathrm{sd}[F_j]\le d_j$. The leg
parameters are fitted on the legs of the development flights. Because the variance of one long leg can exceed that of
the two legs it replaces, removing a task can make a route infeasible under this model, so every route is filtered once
more before publication. With the cost-aware fallback, moment propagation accepts a given-up task if its nominal finish
meets the deadline and $\mathrm{E}[F_j]+z_j\,\mathrm{sd}[F_j]\le d_j$. Online calibration starts from one error model
fitted on both disturbance levels, with the weight of $n_0$ pairs, and refits Eq.~\eqref{eq:errmodel} at every call with
the pairs of the tasks completed so far, so it needs neither the disturbance level nor a level-specific safety factor.
The selected values are $c=0.45$ and 0.6, $z=1$ and 0.5 for moment propagation, and $n_0=50$ with $z=0.5$ for online
calibration, under mid and high disturbances. On development missions, the three static models differed by less than
0.2\,pp. CBBA bids the largest increase of the route value over all feasible insertion positions, plans from scratch at
every call and stops at convergence or at the hard budget. Without the bid warping of Johnson et
al.~\cite{johnson2012nonsubmodular}, it did not converge on some development snapshots. Within \mname{}, the stopping
rule uses a zero mean error, because the effective deadlines already contain it.

\subsection{Tuned competitors of studies 5 and 7}
In study~5, the proportional buffer was tuned with $c\in\{0.1,0.2,0.3,0.45,0.6,0.75,0.9,1.2\}$, with or without the
nominal fallback, and \mname{} without fallback with $z\in\{-0.5,0,0.5,1,1.5\}$, for every $\lambda\in\{0.5,1,2\}$ and
disturbance level. A grid that first ended at $c=0.3$ and $z=0$ was extended after both values were selected on its
edge, and the selection did not change. For $\lambda=1$, the selected buffer was \val{FSelPropOneMid} under mid and
\val{FSelPropOneHigh} under high disturbances, and \mname{} without fallback had \val{FSelCCnfOneMid} and
\val{FSelCCnfOneHigh}. Studies 5 and 6 ran the cost-aware fallback with the proportional cost model,
$\theta_j=\lambda/(1+\lambda)$, but scored the runs with $u_\lambda$, which charges every late arrival a fixed cost;
study~7a removes this mismatch. The fixed-cost rule of study~7a uses the mean reward $\bar w=2.1$ of the generator,
because the total reward of a mission is not known while tasks are still arriving, whereas $u_\lambda$ charges the
realised mean reward $R/n$ of the mission. Over the held-out missions, $R/n$ was \val{UCWnMean} on average, with a
coefficient of variation of \val{UCWnCv}\% and a range of \val{UCWnMin} to \val{UCWnMax}, so study~7a is also reported
with $u'_\lambda=U-\lambda\bar w\,n_{\mathrm{late}}/R$. In study~7a, the development selection of the tuned buffer
differed between $u_\lambda$ and $u^w_\lambda$ only at $\lambda=1$ under high disturbances, with $c=0.45$ for $u_1$ and
$c=0.6$ for $u^w_1$, both without fallback.

Study~7b keeps the error model, the safety factor and every tuned parameter of the calibrated high level. Speed and
service errors follow a Student-$t$ law with three degrees of freedom at the same standard deviation; rare long delays
add, with probability 0.1, an exponential delay with a mean of 90\,s to a leg, as a detour or a holding pattern would; a
persistent speed bias per UAV and mission with a standard deviation of 8\% correlates the errors along a route; and a
Pareto dissemination delay with the median of the lognormal one and a tail index of 1.5 replaces the lognormal delay.
It compares the nominal model, the constant and the proportional buffer, \mname{} with the nominal and with the
cost-aware fallback ($\lambda=1$, both cost models) and moment propagation with the cost-aware fallback.

\subsection{Protocol}
Studies 2, 3, 4, 5 and 6 used the held-out seeds 5000--5029, 7000--7029, 8000--8029, 10000--10029 and 12000--12029,
and studies 7a and 7b used 13000--13029 and 14000--14029. The tier C part of study~2 flew 30 missions with mid
disturbances and a latency factor of 50 (seeds 6000--6009), and study~4C used seeds 9000--9011. The protocols were not
lodged with an external registry. They are available with the code, together with the commit of every freeze and every
deviation. A tier C run in which a vehicle does not take off in time is repeated once, like a run with an invalid
real-time factor. Tier B runs four missions in parallel on an Intel Core i7-11800H laptop, with the methods of a
configuration run in turn. A separately written validator found no infeasible plan among \val{JValidatorPlans}
development plans. Because the stopping rules read a wall clock, the mean utility of a method varied by up to
\val{JRepeatDiffMax}\,pp between two runs of the same missions. Equivalence (H3) was tested with two one-sided tests.

\subsection{Further results of studies 2 to 4}
In study~2, the 90\% interval of the stopping rule minus a fixed 40\,ms was \val{KTostCI}\,pp, inside the
equivalence margin of $\pm0.5$\,pp. \mname{} without fallback served \val{KLateProp} tasks per mission after their
deadline, against \val{KLateForty} for ALNS with 40\,ms. In the tier C part of study~2, the Holm-adjusted $p$ of H4 was
\val{KCPForty}, and about \val{KCNPowerForty} missions per method would have been needed to confirm the difference that
tier B predicted. In study~4, the interval of \mname{} minus the proportional buffer was \val{XCIPropBufAll}\,pp
(Holm-adjusted $p={}$\val{QPPropBufAll}) and that of \mname{} minus moment propagation \val{XCIMPAll}\,pp. The loss of
moment propagation without fallback had an interval of \val{XCIFbMPAll}\,pp. \mname{} ran \val{QItMsCC} search iterations per millisecond and
moment propagation \val{QItMsMP}. Online calibration served \val{XDabsOnlCC}\,pp less than the model with a known
disturbance level, and the learned stopping rule added \val{XDabsFullCC}\,pp. In study~4C, the $p$-value of H14 was
\val{GPNom}. The advantage over the nominal model that sized the study was \val{XDabsTcDevNom}\,pp on development seeds.
Over the full tier B grid, the advantage in this regime was \val{XDabsFortyHighWHalf}\,pp, which would need about
\val{HNPowerNomGrid} missions per method to be confirmed in tier C. The interim analysis did not cross the
O'Brien--Fleming bound. The illustrative flights of Fig.~\ref{fig:demo} ran with the graphical interface, at a
real-time factor of \val{ZDemoRTFNom} to \val{ZDemoRTFProp}.

\subsection{Further results of studies 5 and 6}
In study~5, the advantage of the cost-aware fallback at $\lambda=1$ came from high disturbances, where it served
\val{XDabsFivePropOneHigh} and \val{XDabsFiveCCnfOneHigh}\,pp more than the tuned proportional buffer and the tuned
safety factor. Under mid disturbances, the methods were on par, and with the original windows both tuned competitors
were better, by \val{XDabsFivePropOneMidWOne} and \val{XDabsFiveCCnfOneMidWOne}\,pp. At $\lambda=0.5$, the interval
against the tuned buffer was \val{XCIFivePropHalf}\,pp. At $\lambda=2$, the cost-aware fallback stayed
\val{XDabsFiveCCfbTwo}\,pp above the nominal fallback and matched the tuned safety factor, while the tuned buffer was
better (\val{XCIFivePropTwo}\,pp) and moment propagation without fallback by \val{XDabsFiveMPnfTwo}\,pp. Tuned for
$\lambda=1$, the proportional buffer stayed \val{XDabsFivePropFixHalf}\,pp behind the cost-aware fallback at
$\lambda=0.5$. A wrong late cost also costs the cost-aware fallback, \val{XDabsFiveCAFixTwo}\,pp when it was run with
$\lambda=1$ and scored at $\lambda=2$, and when a late arrival costs nothing, the rule run with $\lambda=1$ served
\val{FUCostCCfb}\,pp less utility than the nominal fallback.

In study~6, CBBA on the nominal model reached almost the utility of ALNS (\val{WCICBBAFortyZero}\,pp), and \mname{}
served \val{WDabsCBBAEventZero}\,pp more than CBBA with the event trigger at $\lambda=0$. CBBA served \val{WLateCBBA}
tasks late per mission and \mname{} \val{WLateCA}. Under high disturbances with halved windows at $\lambda=0$, \mname{}
and CBBA were on par (\val{WDCBBAZeroHighWHalf}\,pp), because a fallback for $\lambda=1$ gives up tasks that a late
arrival would not penalise at $\lambda=0$. Within moment propagation, the cost-aware fallback lost
\val{WDabsMPFbZero}\,pp against the nominal fallback at $\lambda=0$ and gained \val{WDabsMPFbTwo}\,pp at $\lambda=2$.
The interval of the closed form minus moment propagation, both with the cost-aware fallback, was \val{WCICAMPCAOne}\,pp
at $\lambda=1$, and moment propagation ran \val{WItersMPCA} search iterations per call against \val{WItersCA}.

\subsection{Further results of study 7}
In study~7a, the fixed-cost rule scored \val{USCAFeeOneFee}\% with $u_1$, the proportional rule \val{USCAPropOneFee}\%
and the tuned buffer \val{USPropOneFee}\%; the Holm-adjusted $p$ of H24 was \val{UPHTwentyFour}. H26 had a
Holm-adjusted $p$ of \val{UPHTwentySix} for non-inferiority and $p={}$\val{UPSupHTwentySix} for superiority. The
fixed-cost rule scored \val{UDFeeCApropTwoFee}\,pp above the proportional rule with $u_2$, and the proportional rule
scored \val{UDabsFeeCApropOneW} and \val{UDabsFeeCApropTwoW}\,pp above the fixed-cost rule with $u^w_1$ and $u^w_2$.
With $u_2$, the fixed-cost rule minus the tuned buffer was \val{UDFeePropTwoFee}\,pp, whereas the proportional rule
stayed behind it (\val{UDPropPropTwoFee}\,pp). With $u^w_2$, the proportional rule stayed \val{UDabsPropPropTwoW}\,pp
behind the buffer tuned for that score. With $u'_1$, the differences of H24 and H25 were \val{UCDCApropOne} and
\val{UCDPropOne}\,pp. In study~7b, the nominal model served \val{UBUFortyGauss}\% under the calibrated errors and
between \val{UBUFortyOutlier}\% with rare long delays and \val{UBUFortyTthree}\% with Student-$t$ errors. The cost-aware
fallback scored \val{UDabsBCAPropTthree} to \val{UDabsBCAPropOutlier}\,pp more than the proportional buffer of study~4,
and rare long delays raised the mean finish-time error of \mname{} from \val{UBErrGauss} to \val{UBErrOutlier}\,s.
